% titel: Blatt 10, A8: pH von 50 mg CH3COOH + 10 mg NaAcetat
% Teil II: der durchgerechnete Weg.

\subsection*{Gegeben und gesucht}
50 mg CH$_3$COOH und 10 mg CH$_3$COONa, mit Wasser auf 100 mL. Gesucht pH.

\subsection*{Rechnung}
$M(\mathrm{CH_3COOH})\approx60.05$, $M(\mathrm{CH_3COONa})\approx82.03$
g/mol; $pK_s=4.75$.
\begin{align*}
  n(\mathrm{HAc}) &= \frac{0.050}{60.05} = 8.33\cdot10^{-4}\ mol,\\
  n(\mathrm{Ac^-}) &= \frac{0.010}{82.03} = 1.22\cdot10^{-4}\ mol.
\end{align*}
Puffergleichung (V kürzt sich):
\[
  pH = pK_s + \lg\frac{n(\mathrm{Ac^-})}{n(\mathrm{HAc})}
     = 4.75 + \lg\frac{1.22\cdot10^{-4}}{8.33\cdot10^{-4}}
     = 4.75 + \lg 0.146 .
\]
\[
  pH \approx 4.75 - 0.835 \approx 3.92 .
\]

\begin{ergebnis}
$pH \approx 3.9$ (Essigsäure $\gg$ Acetat, daher im saueren Bereich).
\end{ergebnis}

\begin{kontrolle}
Die Säure ist mehrfach im Überschuss gegenüber dem Acetat
(Ratio$\approx0.146$), also liegt der pH deutlich unter $pK_s(=4.75)$.
$pH\approx3.9$ ist plausibel.
\end{kontrolle}

\begin{merksatz}
Mengen der Pufferkomponenten aus $n=m/M$; V kürzt sich. $pH=pK_s+\lg
\frac{n(\text{Base})}{n(\text{Säure})}$.
\end{merksatz}
